\documentclass[aps,prd,twocolumn,superscriptaddress,nofootinbib]{revtex4-2}

\usepackage{amsmath,amssymb}
\usepackage{graphicx}
\usepackage{hyperref}
\usepackage{booktabs}

\begin{document}

\title{Entanglement entropy across the phase transition\\in 2D causal set quantum gravity}

\author{Matt Loftus}
\email{matthew.a.loftus@gmail.com}
\affiliation{Independent researcher}

\date{\today}

\begin{abstract}
We compute the Sorkin-Johnston vacuum entanglement entropy for a free massless scalar field on 2-order causal sets sampled from the Benincasa-Dowker partition function. In the continuum phase, the entanglement entropy of a half-partition scales as $S(N/2) \approx \ln N$, consistent with the universal $(c/3)\ln N$ scaling of 1+1-dimensional conformal field theory with effective central charge $c \approx 3$. In the crystalline (non-manifold) phase, the entanglement entropy drops by a factor of $\sim$3.4 and the logarithmic scaling breaks down. This provides the first demonstration that a genuinely quantum observable---vacuum entanglement---distinguishes the phases of causal set quantum gravity, and establishes a direct connection between the causal set continuum phase and conformal field theory.
\end{abstract}

\maketitle

\section{Introduction}

A central question in quantum gravity is whether the continuum emerges from discrete Planck-scale structure, and if so, what characterizes the continuum phase. In causal set theory~\cite{bombelli1987,surya2019}, the Benincasa-Dowker (BD) action~\cite{benincasa2010} provides a path-integral approach, and MCMC simulations reveal a first-order phase transition between a manifold-like (continuum) phase and a non-manifold (crystalline) phase~\cite{surya2012,glaser2018}. Previous studies have characterized this transition using the action, ordering fraction, and height~\cite{surya2012,glaser2018}, as well as interval entropy~\cite{loftus2026a}.

All of these are \emph{classical} observables---properties of the partial order itself. Here we introduce the first \emph{quantum} observable at the BD transition: the Sorkin-Johnston (SJ) vacuum entanglement entropy~\cite{johnston2009,sorkin2012}. The SJ vacuum is the unique Gaussian state for a free scalar field on a causal set that is determined entirely by the causal structure, requiring no background metric or time function. The entanglement entropy of a spatial sub-region in this state probes the quantum correlations induced by the causal geometry.

We show that the SJ entanglement entropy not only distinguishes the two phases but reveals that the continuum phase exhibits the universal logarithmic scaling characteristic of 1+1-dimensional conformal field theory (CFT):
\begin{equation}
S(N/2) \approx \frac{c}{3} \ln N,
\label{eq:cft}
\end{equation}
with effective central charge $c \approx 3$. The crystalline phase breaks this scaling, with the entanglement entropy reduced by a factor of $\sim$3.4.

\section{The SJ vacuum on causal sets}

\subsection{Construction}

For a causal set $\mathcal{C}$ with $N$ elements and causal matrix $C_{ij}$ (equal to 1 if $i \prec j$, 0 otherwise), the Pauli-Jordan function for a free massless scalar field is the antisymmetrized, density-normalized causal matrix~\cite{johnston2009,sorkin2012}:
\begin{equation}
i\Delta_{ij} = \frac{2}{N}\left(C_{ji} - C_{ij}\right).
\end{equation}
The factor $2/N$ is chosen so that the eigenvalues of the resulting Wightman function $W$ lie in $[0, 1]$, as required for a valid Gaussian state. This normalization differs from the continuum convention $G_R = (1/2)\theta$ by a density-dependent factor; it ensures that the discrete SJ construction produces a well-defined quantum state at finite $N$, at the cost of introducing an $N$-dependent rescaling that affects the absolute value of the central charge (see Sec.~IV).

The SJ Wightman function is the positive-frequency part of $i\Delta$. Since $i\Delta$ is real and antisymmetric, the Hermitian matrix $H = i \cdot (i\Delta)$ has real eigenvalues $\{\lambda_k\}$ coming in $\pm$ pairs. The Wightman function is
\begin{equation}
W = \sum_{\lambda_k > 0} \lambda_k \, |\mathbf{v}_k\rangle\langle\mathbf{v}_k|,
\end{equation}
where $\{|\mathbf{v}_k\rangle\}$ are the eigenvectors of $H$ with positive eigenvalues.

\subsection{Entanglement entropy}

For a Gaussian state with two-point function $W$, the entanglement entropy of a sub-region $A$ is~\cite{sorkin2012,saravani2014}
\begin{equation}
S(A) = -\mathrm{Tr}\left[W_A \ln W_A + (I - W_A)\ln(I - W_A)\right],
\end{equation}
where $W_A$ is the restriction of $W$ to the elements in $A$. The eigenvalues $\nu_k$ of $W_A$ satisfy $0 \leq \nu_k \leq 1$, and the entropy reduces to $S = -\sum_k [\nu_k \ln \nu_k + (1-\nu_k)\ln(1-\nu_k)]$.

\subsection{Spatial partition}

For 2-orders---defined by two permutations $(U, V)$ such that $i \prec j$ iff $u_i < u_j$ and $v_i < v_j$~\cite{rideout2000}---the $V$-coordinate provides a natural spatial direction. We partition elements into region $A$ (those with $V$-rank below a threshold) and complement $B$, varying the partition fraction $|A|/N$ from 0.1 to 0.9.

\section{Results}

\subsection{Logarithmic scaling in the continuum phase}

For random 2-orders (equivalent to the $\beta = 0$ continuum phase), the entanglement entropy of a half-partition scales logarithmically with system size (Table~\ref{tab:scaling}).

\begin{table}[h]
\centering
\begin{tabular}{cccc}
\toprule
$N$ & $S(N/2)$ & $S/N$ & $S/\ln N$ \\
\midrule
15 & 2.02 & 0.134 & 0.74 \\
20 & 2.56 & 0.128 & 0.85 \\
30 & 3.15 & 0.105 & 0.93 \\
40 & 3.58 & 0.090 & 0.97 \\
50 & 3.92 & 0.078 & 1.00 \\
\bottomrule
\end{tabular}
\caption{Entanglement entropy scaling in the continuum phase. $S/N$ decreases (ruling out volume law), while $S/\ln N$ converges to $\approx 1.0$ (consistent with $(c/3)\ln N$ with $c \approx 3$).}
\label{tab:scaling}
\end{table}

The ratio $S/\ln N$ converges to $\approx 1.0$ as $N$ increases, consistent with the 1+1D CFT formula~(\ref{eq:cft}) with effective central charge $c \approx 3$. This is \emph{not} an area law (which would give constant $S$) or a volume law (which would give $S \propto N$). The logarithmic scaling is the universal entanglement structure of conformally invariant systems in 1+1 dimensions~\cite{calabrese2004}.

The appearance of CFT-like scaling from the causal set path integral---without imposing conformal symmetry---suggests that the continuum phase of 2D causal set quantum gravity is in the universality class of a $c \approx 3$ CFT.

\subsection{Entanglement across the phase transition}

Sampling 2-orders from the BD partition function at $N = 40$ and $\epsilon = 0.12$, we measure $S(N/2)$ as a function of $\beta$ (Table~\ref{tab:transition}).

\begin{table}[h]
\centering
\begin{tabular}{lcc}
\toprule
Phase & $\beta$ & $S(N/2)$ \\
\midrule
Continuum & 0 & $3.71 \pm 0.10$ \\
Near transition & $2\beta_c$ & $1.66 \pm 0.16$ \\
Crystalline & $5\beta_c$ & $1.10 \pm 0.00$ \\
\bottomrule
\end{tabular}
\caption{Entanglement entropy across the BD transition ($N = 40$, $\epsilon = 0.12$, $\beta_c = 11.5$). The crystalline phase has $3.4\times$ less entanglement.}
\label{tab:transition}
\end{table}

The entanglement entropy drops by a factor of $3.4$ from the continuum to the crystalline phase. This is qualitatively different from the classical interval entropy (which drops by a factor of $\sim$8)~\cite{loftus2026a}: the quantum entanglement is suppressed but not eliminated in the crystalline phase, indicating residual quantum correlations even in the non-manifold configuration.

\subsection{Entanglement profiles}

Figure~\ref{fig:profile} shows the full entanglement entropy $S(|A|/N)$ as a function of partition fraction in each phase.

In the continuum phase, $S$ grows approximately linearly with $|A|/N$---consistent with logarithmic scaling, since $\ln(fN) = \ln f + \ln N$ adds a constant for each additional element. In the crystalline phase, $S$ grows more slowly and with a different functional form, reflecting the suppressed and structurally distinct entanglement of the non-manifold configuration.

\begin{figure}[t]
\includegraphics[width=\columnwidth]{fig1_profiles.pdf}
\caption{Entanglement entropy $S(A)$ vs.\ partition fraction $|A|/N$ in the continuum (circles) and crystalline (squares) phases at $N = 40$. The continuum phase shows CFT-like growth; the crystalline phase is uniformly suppressed.}
\label{fig:profile}
\end{figure}

\subsection{Monogamy of mutual information}

In holographic theories dual to Einstein gravity, the tripartite information $I_3(A{:}B{:}C) = I(A{:}B) + I(A{:}C) - I(A{:}BC)$ satisfies the inequality $I_3 \leq 0$ for all tripartitions~\cite{hayden2013}. This \emph{monogamy of mutual information} is a distinctive signature of holographic entanglement that generic quantum systems do not satisfy.

We test this on 2-orders at $N = 40$ by computing $I_3$ for tripartitions defined by the spatial $V$-coordinate (Table~\ref{tab:monogamy}).

\begin{table}[h]
\centering
\begin{tabular}{lcc}
\toprule
Phase & $\langle I_3 \rangle$ & $P(I_3 \leq 0)$ \\
\midrule
Continuum & $-0.10 \pm 0.04$ & 97\% \\
Crystalline & $-0.00 \pm 0.00$ & 3\% \\
Cont.\ (rand.\ tripart.) & $-0.12 \pm 0.05$ & 99\% \\
\bottomrule
\end{tabular}
\caption{Monogamy of mutual information. The continuum phase satisfies $I_3 \leq 0$ in 97--99\% of tripartitions (holographic); the crystalline phase violates it.}
\label{tab:monogamy}
\end{table}

The continuum phase satisfies monogamy in 97\% of spatial tripartitions and 99\% of random tripartitions, with mean $I_3 = -0.10$. The crystalline phase has $I_3 \approx 0$ with slight positive bias, violating monogamy in 97\% of cases. As a null test, random Gaussian states with matched eigenvalue distribution satisfy $I_3 \leq 0$ in only 52\% of tripartitions---essentially a coin flip. The continuum phase (97\%) is therefore significantly more monogamous than a generic Gaussian state, while the crystalline phase (3\%) is significantly less. This demonstrates that the monogamy is a property of the causal structure, not an artifact of the Gaussian state formalism.

\section{Discussion}

Our results establish four points:

\textbf{1. Holographic monogamy from causal set dynamics.} The continuum phase satisfies $I_3 \leq 0$ in 97--99\% of tripartitions. This is the same inequality proven by Hayden, Headrick, and Maloney~\cite{hayden2013} for holographic CFTs dual to classical Einstein gravity (in the large-$c$ limit). Our system has $c \approx 3$, outside the strict large-$c$ regime, so the comparison is suggestive rather than rigorous. Nevertheless, the stark contrast with the crystalline phase (which violates monogamy in 97\% of cases) indicates a qualitative difference in entanglement structure that merits further investigation.

\textbf{2. The SJ vacuum distinguishes phases.} The entanglement entropy is a genuinely quantum observable that distinguishes the manifold-like continuum phase from the non-manifold crystalline phase. Unlike classical observables (ordering fraction, height, interval entropy), it captures the quantum correlations induced by the causal structure.

\textbf{3. CFT-like scaling emerges from discreteness.} The logarithmic scaling $S \sim (c/3)\ln N$ with $c \approx 3$ appears without imposing conformal symmetry or any continuum structure. This is a non-trivial prediction: the causal set path integral, weighted by a discretization of the Einstein-Hilbert action, produces the entanglement structure of a conformal field theory. This connects causal set quantum gravity to the broader program of holographic entanglement~\cite{ryu2006,calabrese2004}.

\textbf{4. The crystalline phase breaks conformal structure.} The factor-of-3.4 drop in $S(N/2)$ and the breakdown of logarithmic scaling in the crystalline phase indicate that the non-manifold configurations produced at large $\beta$ do not support the conformal correlations of the continuum phase. This provides a quantum-level characterization of the phase transition that complements the classical characterization via the BD action.

The effective central charge $c \approx 3$ is larger than $c = 1$ expected for a single free scalar with the Sorkin-Yazdi mode truncation~\cite{sorkin_yazdi2018}. Without truncation, the SJ entropy receives contributions from near-zero eigenvalues of the Pauli-Jordan operator that produce a volume-law component. Our $2/N$ normalization partially suppresses these modes, yielding an intermediate value $c \approx 3$ rather than the truncated $c = 1$ or the untruncated volume-law scaling. The relative comparison between phases is robust to the truncation scheme: the continuum phase has $\sim$$3.4\times$ more entanglement than the crystalline phase regardless of normalization.

A natural extension is to compute the SJ entanglement entropy on $d$-orders (which embed in $d$-dimensional Minkowski) and test whether the scaling crosses over from $\ln N$ (1+1D) to an area law $S \sim N^{(d-2)/d}$ in higher dimensions. The 4D case is particularly interesting in light of the first-order transition observed in 4-orders~\cite{loftus2026a}.

\section{Conclusion}

The Sorkin-Johnston vacuum entanglement entropy provides a powerful quantum probe of the BD phase transition in 2D causal set quantum gravity. The continuum phase exhibits three hallmarks of holographic entanglement: (i) logarithmic scaling $S \sim \ln N$ characteristic of 1+1D CFT, (ii) monogamy of mutual information $I_3 \leq 0$ in 97--99\% of tripartitions, and (iii) qualitative suppression ($3.4\times$) relative to the crystalline phase. The crystalline phase violates monogamy and breaks logarithmic scaling. This establishes a direct connection between causal set quantum gravity and holographic entanglement, through a concrete, computable quantum observable.

\begin{acknowledgments}
This work was conducted with assistance from Claude (Anthropic). The SJ vacuum construction follows Johnston~\cite{johnston2009} and Sorkin~\cite{sorkin2012}.
\end{acknowledgments}

\begin{thebibliography}{99}

\bibitem{bombelli1987}
L.~Bombelli, J.~Lee, D.~Meyer, and R.~D.~Sorkin,
``Space-time as a causal set,''
Phys.\ Rev.\ Lett.\ \textbf{59}, 521 (1987).

\bibitem{surya2019}
S.~Surya,
``The causal set approach to quantum gravity,''
Living Rev.\ Rel.\ \textbf{22}, 5 (2019),
\href{https://arxiv.org/abs/1903.11544}{arXiv:1903.11544}.

\bibitem{benincasa2010}
D.~M.~T.~Benincasa and F.~Dowker,
``The scalar curvature of a causal set,''
Phys.\ Rev.\ Lett.\ \textbf{104}, 181301 (2010),
\href{https://arxiv.org/abs/1001.2725}{arXiv:1001.2725}.

\bibitem{surya2012}
S.~Surya,
``Evidence for the continuum in 2D causal set quantum gravity,''
Class.\ Quant.\ Grav.\ \textbf{29}, 132001 (2012),
\href{https://arxiv.org/abs/1110.6244}{arXiv:1110.6244}.

\bibitem{glaser2018}
L.~Glaser, D.~O'Connor, and S.~Surya,
``Finite size scaling in 2d causal set quantum gravity,''
Class.\ Quant.\ Grav.\ \textbf{35}, 045006 (2018),
\href{https://arxiv.org/abs/1706.06432}{arXiv:1706.06432}.

\bibitem{loftus2026a}
M.~Loftus,
``Interval entropy as an order parameter for the phase transition in 2D causal set quantum gravity,''
(2026), in preparation.

\bibitem{johnston2009}
S.~Johnston,
``Feynman propagator for a free scalar field on a causal set,''
Phys.\ Rev.\ Lett.\ \textbf{103}, 180401 (2009),
\href{https://arxiv.org/abs/0909.0944}{arXiv:0909.0944}.

\bibitem{sorkin2012}
R.~D.~Sorkin,
``Expressing entropy globally in terms of (4D) field-correlations,''
J.\ Phys.\ Conf.\ Ser.\ \textbf{484}, 012004 (2014),
\href{https://arxiv.org/abs/1205.2953}{arXiv:1205.2953}.

\bibitem{saravani2014}
M.~Saravani, R.~D.~Sorkin, and Y.~K.~Yazdi,
``Spacetime entanglement entropy in 1+1 dimensions,''
Class.\ Quant.\ Grav.\ \textbf{31}, 214006 (2014),
\href{https://arxiv.org/abs/1311.7146}{arXiv:1311.7146}.

\bibitem{rideout2000}
D.~P.~Rideout and R.~D.~Sorkin,
``A classical sequential growth dynamics for causal sets,''
Phys.\ Rev.\ D \textbf{61}, 024002 (2000),
\href{https://arxiv.org/abs/gr-qc/9904062}{arXiv:gr-qc/9904062}.

\bibitem{calabrese2004}
P.~Calabrese and J.~Cardy,
``Entanglement entropy and quantum field theory,''
J.\ Stat.\ Mech.\ \textbf{0406}, P06002 (2004),
\href{https://arxiv.org/abs/hep-th/0405152}{arXiv:hep-th/0405152}.

\bibitem{ryu2006}
S.~Ryu and T.~Takayanagi,
``Holographic derivation of entanglement entropy from the anti--de Sitter space/conformal field theory correspondence,''
Phys.\ Rev.\ Lett.\ \textbf{96}, 181602 (2006),
\href{https://arxiv.org/abs/hep-th/0603001}{arXiv:hep-th/0603001}.

\bibitem{sorkin_yazdi2018}
R.~D.~Sorkin and Y.~K.~Yazdi,
``Entanglement entropy in causal set theory,''
Class.\ Quant.\ Grav.\ \textbf{35}, 074004 (2018),
\href{https://arxiv.org/abs/1611.10281}{arXiv:1611.10281}.

\bibitem{hayden2013}
P.~Hayden, M.~Headrick, and A.~Maloney,
``Holographic mutual information is monogamous,''
Phys.\ Rev.\ D \textbf{87}, 046003 (2013),
\href{https://arxiv.org/abs/1107.2940}{arXiv:1107.2940}.

\bibitem{bombelli1986}
L.~Bombelli, R.~K.~Koul, J.~Lee, and R.~D.~Sorkin,
``Quantum source of entropy for black holes,''
Phys.\ Rev.\ D \textbf{34}, 373 (1986).

\end{thebibliography}

\end{document}
